% !TeX root = ../papers/nc-p3.tex
\section{一辺だけからの二元障害の伝播と指数成長}
\label{sec:single-binary-edge-obstruction}
\begin{theorem}\label{thm:single-binary-edge-obstruction}
標準有限Frobenius四窓の一辺に二元形式の商があるとする。
その辺の六関係は商のHankel六空間に同型に写り、三次形式の四文字の側で続く隣の辺には、
その有理正規三次曲線の点を因子とする純関係がないとする（逆向きにも同じ主張が成り立つ）。
このとき元の窓は指数成長し、標準AS正則にはならない。
特に全辺の点グラフが有限なFrobenius窓では、\emph{一辺だけ}の二元商の存在で非正則性が決まる。
長さ $4m+1$ のある始点で経路次元は $4\cdot3^m$ 以上で、長さ29ですでに標準次元を超える。
有限グラフの場合、もとの四窓は三次形式の側の反対二位置で必ず交換対称な表示を持つ。
ここでは残る三辺の商、正則性、交換対称性を事前に仮定しない。
\end{theorem}
\begin{proof}[ベクトル値二元逆系の分裂]
向きを $U,Q,V,W$ と書き、点座標側の $U$ に $P_2$ の三平面を取り、$Q=P_3$ とする。
関係の空間はこれらの双対に入る。指定の商の双対は $P_2\otimes P_3\to P_5$ の乗法である。
$U$ を $P_2$ に制限した四テンソルの値は $p,q$ に関して $pq$ だけに依存する。
続く辺の六関係を $R_1\subset P_3^*\otimes V^*$ とし、その零化空間 $B\subset P_3\otimes V$ を取る。
仮定により次数三の評価 $B\otimes\mathcal O\to V\otimes\mathcal O(3)$ は
$\mathbb P^1$ 上全射である。ある点 $t$ で全射でなければちょうど
$\operatorname{ev}_t\otimes v^*\in R_1$ ($v^*\ne0$) となるためである。
この十列から四行への全射の核束を $\bigoplus_{j=1}^6\mathcal O(-d_j)$ と分裂させる。
十列は独立なので $d_j\ge1$、次数から $\sum d_j=12$ である。

$n=3+j$ とし、$M_n\subset P_n^*\otimes V^*$ を $B P_j$ の零化空間とする。
$P_j$ の因子との全ての縮約が $R_1$ に入ることと同じである。
上の束列のコホモロジーから
\[
       M_{3+j}^*=H^1\!\left(\bigoplus\mathcal O(j-d_i)\right),
       \qquad \dim M_{3+j}=\sum_i\max(d_i-j-1,0).
\]
圧縮した四テンソルの第四文字は $M_5$ に単射で写る。
実際 $P_2$ との全縮約は $R_1$ に入り、最初の二位置の六関係を単射に圧縮したので
反対二位置の六関係も第四文字の全支持も失われない。
従って $\sum(d_i-3)_+\ge4$。六つの $d_i-1$ の和が6であるから、可能なのは
$(d_i)=(1,1,1,1,1,7)$ だけである。
この唯一の長い直和因子とSerre双対性により $M_8$ は一本のベクトル値八次汎関数
$\Phi\in P_8^*\otimes V^*$ で張られ、$3\le n\le8$ について
$P_{8-n}\to M_n$, $h\mapsto\Phi\mathbin{\lrcorner}h$ は同型である。
実際積による遷移は束列の接続写像と可換で、長い直和因子の双対では通常の二元乗法である。

これで第四位置は $M_5\simeq P_3$ と同定され、各 $v\in V$ を $D=\Phi(v)\in P_8^*$ として
圧縮した窓は $D(pqq')$ となる。$V\to P_8^*$ は単射である。核があれば
$R_1=M_3$ の一文字支持が落ちる。像を $S$ とする。
$R_1$ とその次の辺 $R_2$ は、$Q,S$ と $S,Q$ における同じ写像
\[
           (q,D)\longmapsto D\mathbin{\lrcorner}q\in P_5^*
\]
の双対像とその転置である。$R_1$ については $M_3$ の表示、$R_2$ については
対になる最初の辺の圧縮が単射であることによる。
$R_3$ を $P_3^*\otimes P_2^*$ へ圧縮した像も二元乗法の双対像になる。
実際、圧縮テンソルのその側の像は $R_1$ との元の完全対合から圧縮 $R_3$ の像全体で、
圧縮窓の式と $P_5\simeq M_3$ の単射性から六次元のHankel空間に等しい。
\end{proof}
\begin{proof}[切り離せる三文字ブロック]
$U$ 位置の生成文字を三つに割った経路代数を考える。
この位置と両隣の $P_3$ の間の二元乗法の双対関係を消すため、双対の経路空間では
各隣接組の対角式 $\Delta(u,q)$ を掛ける。内部の $u$ の次数2は両側の対角式でちょうど尽き、
各 $q$ では次数2が残る。$S$ 位置を挟む三文字に
$F\in P_2\otimes S\otimes P_2$ を置く。
両隣の対角式を掛けても中央の両辺の関係を消す十分条件は
\[
 \bar\gamma:P_2\otimes S\to P_6^*,\qquad (r,D)\longmapsto D\mathbin{\lrcorner}r
\]
をどちらの端から作用させても零となることである。実際、対角式の係数に現れる
二つの一次式を $P_5$ に掛けたものは $P_6$ を張る。
両端について対称な $F$ の空間は $6\cdot4=24$ 次元で、一方からの収縮先は
$7\cdot3=21$ 次元、他方の条件は対称性から従う。よって少なくとも三つの独立な $F$ がある。

長さ $4m+1$ の $u,q,S,q',u,\ldots,q,S,q',u$ という並びで $m$ 個のブロックを
独立に選び、全対角式を掛ける。両端の $u$ には各々まだ任意の一次式を掛けられる。
異なるブロックは変数を共有せず、固定した対角式の積による乗法も単射なので、
次元は少なくとも $4\cdot3^m$。$m=7$ で $8748>\binom{32}3=4960$ となる。
元の窓はこの経路代数に全射するので同じ下界を持つ。
\end{proof}
\begin{proof}[有限グラフでの対称性]
有限グラフの全関係に純テンソルがないので定理の二つの仮定は自動的である。
$Q,S,Q$ の両辺が互いに転置であることから、その四次元の三次交叉は両端の交換で保たれる。
$P_2\subset U$ の三文字による四次窓の縮約は端の交換で不変で互いに独立。
残りの一文字を選び直せば、全四文字が不変か、第四文字だけ符号が変わるかのいずれかである。
後者なら $Q$ 両端の交換と $U$ でのこの射影反射、$S$ での恒等の組が窓を保つ。
これは四辺の点列を逆順にするが、$S$ 位置への射影は点列の閉埋込みで、そこで各点を
スキームとして固定するため、$U$ 位置の配置も点ごとに固定しなければならない。
配置への二次式制限は単射であり、全ての点で恒等と比例する行列はスカラーなので
この反射は不可能である。従って元の窓自身で $Q$ の二位置は交換対称となる。
\path{computations/verify_single_binary_edge_obstruction.py}は分裂次数、対角式の両係数と成長下界を検算する。
\end{proof}
